\subsubsection{跨字边界的确定性映射}
以 tile 起点的前 8 个元素 $y_0$ 到 $y_7$ 为例，每个字段使用其 12 位二进制补码。第一个总线字容纳 $y_0$、$y_1$ 和 $y_2$ 的低 8 位；第二个字包含 $y_2$ 的高 4 位、$y_3$、$y_4$ 和 $y_5$ 的低 4 位；第三个字包含 $y_5$ 的高 8 位、$y_6$ 与 $y_7$。下标决定所有边界，与输出值无关。

\begin{figure}[H]\centering
\begin{tikzpicture}[x=.36cm,y=1cm]
\foreach \row in {0,1,2}{\draw[ink](0,-\row*1.1)rectangle(32,-\row*1.1+.65);\node[anchor=east,note]at(-.7,-\row*1.1+.32){字 \row};}
\foreach \a/\b/\lab in {0/12/{$y_0$},12/24/{$y_1$},24/32/{$y_2[7:0]$}}{\draw[fill=pale,draw=ink](\a,0)rectangle(\b,.65);\node[note]at({(\a+\b)/2},.32){\lab};}
\foreach \a/\b/\lab in {0/4/{$y_2$},4/16/{$y_3$},16/28/{$y_4$},28/32/{$y_5$}}{\draw[fill=accent!15,draw=ink](\a,-1.1)rectangle(\b,-.45);\node[note]at({(\a+\b)/2},-.78){\lab};}
\foreach \a/\b/\lab in {0/8/{$y_5[11:4]$},8/20/{$y_6$},20/32/{$y_7$}}{\draw[fill=pale,draw=ink](\a,-2.2)rectangle(\b,-1.55);\node[note]at({(\a+\b)/2},-1.88){\lab};}
\node[note]at(16,-2.85){从左向右表示位号由低到高；字 1 两端分别为 $y_2[11:8]$、$y_5[3:0]$};
\end{tikzpicture}
\caption{8 个定长输出字段在 3 个 32 位字中的精确分布。字编号相对于当前 96 位组起点。}
\end{figure}

最后一个有效位置为 $(i,j)=(98,255)$，对应 $t=15,c=15$，tile 内序号 1583。其位偏移为 18996，起始字为全输出区域的字 9520、字内偏移 20；该元素占该字最高 12 位，恰好结束在分配区域末端。这个索引校验说明 9521 字分配包含全部头部与数据，没有额外尾字，也没有遗漏最后一行或最后一列。

\subsubsection{打包器缓冲与完成条件}
后处理每次交付 4 个数，接口上每个数以 13 位承载，其中最高位须与第 11 位符号一致。打包器先核对这个条件，再提取 4 个 12 位字段形成 48 位有效数据。64 位缓冲在输出一个 32 位字后可剩余 0 位或 16 位；只有剩余空间足够时才接受下一组 48 位。该组织将移位选择限定在少量情形，输出流长度保持固定。

\begin{table}[H]\centering\small
\caption{固定输出的数量关系与完成判定}
\begin{tabularx}{\linewidth}{p{42mm}p{32mm}Y}\toprule
计数对象 & 每 tile & 全任务\\\midrule
矩阵元素 & 1584 & 25344\\
每次 4 元素的输入组 & 396 & 6336\\
32 位数据写出字 & 594 & 9504\\
固定长度字段 & 1 & 16\\
地址分配字数 & -- & 17 字头部加 9504 字数据\\
实际写命令数 & -- & 9521 字另加一次标志失效写，共 9522\\\bottomrule
\end{tabularx}
\end{table}

某 tile 的末字成立需要已接收 396 组元素且缓冲剩余 32 位，不能只观察输入 last 标志提前结束。任务起始先写无效标志，依次输出固定长度字段和完整矩阵，最终提交有效标志；DMA 等待该标志写响应后才发布完成。中途错误使本任务停止继续生成有效完成状态，避免接收方把部分矩阵当作完整结果。
